\section{Placement}
\label{sec:placement}

Once the chunks are defined, the model must decide on which pallet row each
chunk is dropped and where it begins horizontally. Two decision variables
encode this choice:

\vspace{0.5cm}
\begin{lstlisting}
array[1..n_chunks] of var 1..boxes_along_y: y_row;
array[1..n_chunks] of var 1..boxes_along_x: x_start;
\end{lstlisting}
\vspace{0.5cm}

The variable \texttt{y\_row[g]} is the pallet row assigned to chunk \(g\),
and \texttt{x\_start[g]} is the column at which the chunk begins.

An initial formulation may rely on two geometric constraints. The first
ensures that each chunk fits horizontally inside the pallet:

\vspace{0.5cm}
\begin{lstlisting}
constraint forall(g in 1..n_chunks)(
  x_start[g] + chunk_len[g] - 1 <= boxes_along_x
);
\end{lstlisting}
\vspace{0.5cm}

The second forbids two chunks from overlapping on the pallet grid. Each
chunk occupies a horizontal strip of length \texttt{chunk\_len[g]} starting
at column \texttt{x\_start[g]} on row \texttt{y\_row[g]}, which makes every
chunk an axis-aligned rectangle of height \(1\). The global constraint
\texttt{diffn(x, y, dx, dy)} expresses the non-overlap of a collection of
such rectangles: each rectangle is described by an origin
\((\texttt{x[i]}, \texttt{y[i]})\) and dimensions
\((\texttt{dx[i]}, \texttt{dy[i]})\), and \texttt{diffn} enforces that no
two rectangles intersect on the grid. With heights fixed to \(1\) for every
chunk, the constraint expresses the no-overlap requirement directly:

\vspace{0.5cm}
\begin{lstlisting}
include "globals.mzn";   % for diffn

constraint diffn(x_start, y_row, chunk_len, [1 | g in 1..n_chunks]);
\end{lstlisting}
\vspace{0.5cm}

Together with the chunking partition, which forces the chunks to cover all
\texttt{n\_boxes} boxes, these two constraints already describe a correct
placement: the chunks occupy disjoint cells of the pallet, and their
combined length matches the pallet capacity, so every cell ends up covered
exactly once.

However, the structural property that each row receives exactly
\texttt{chunks\_per\_row} chunks is left implicit. It is provable from the
parameter identity \(\texttt{n\_chunks} = \texttt{boxes\_along\_y} \cdot
\texttt{chunks\_per\_row}\) together with the constraints above, but the
solver can only derive it through search.

To expose this property to the propagator from the start, we can refine 
the model by adding a global cardinality constraint:

\vspace{0.5cm}
\begin{lstlisting}
constraint global_cardinality(
  y_row,
  [r | r in 1..boxes_along_y],
  [chunks_per_row | r in 1..boxes_along_y]
);
\end{lstlisting}
\vspace{0.5cm}

The global constraint \texttt{global\_cardinality(arr, vals, counts)}
enforces that in \texttt{arr} each value \texttt{vals[i]} appears exactly
\texttt{counts[i]} times. Here, every row from \(1\) to
\texttt{boxes\_along\_y} must appear exactly \texttt{chunks\_per\_row}
times in \texttt{y\_row}. As soon as some rows fill up during search, the
remaining chunks are immediately forced toward the rows that still have
free slots, which substantially shrinks the branching effort.

\section{Release and Plate Collision}
\label{sec:release}

When a chunk is released, the gripper opens one of its two plates. The
opening side sweeps the adjacent pallet row: if that row already contains a
box overlapping the chunk along \(x\), the plate collides with it. A binary
variable encodes the opening side, with \(0\) meaning the plate opens
upward (toward row \(r-1\)) and \(1\) downward (toward row \(r+1\)):

\vspace{0.5cm}
\begin{lstlisting}
array[1..n_chunks] of var 0..1: plate;
\end{lstlisting}
\vspace{0.5cm}

\subsection{Initial formulation}

The initial formulation expresses the no-collision rule as a single reified
implication over all ordered pairs of distinct chunks:

\vspace{0.5cm}
\begin{lstlisting}
constraint forall(g, g2 in 1..n_chunks where g != g2)(
  (
    ( (plate[g] = 0 /\ y_row[g2] = y_row[g] - 1) \/
      (plate[g] = 1 /\ y_row[g2] = y_row[g] + 1) )
    /\
    x_start[g2] < x_start[g] + chunk_len[g] /\
    x_start[g2] + chunk_len[g2] > x_start[g]
  )
  -> g2 > g
);
\end{lstlisting}
\vspace{0.5cm}

The antecedent describes a collision configuration: chunk \(g_2\) lies in
the row swept by the plate of \(g\), and the two chunks overlap along
\(x\). The first disjunction handles the two opening directions; the two
inequalities express the standard interval overlap test on the
\(x\)-axis. When the antecedent holds, the consequent forces
\(g_2 > g\): the collision can only be tolerated if \(g_2\) is released
\emph{after} \(g\) (so that \(g_2\) is not yet on the pallet at the
moment \(g\) is dropped). The contrapositive captures the physical content:
no chunk released before \(g\) can sit in the target row of \(g\) with
horizontal overlap.

Chunks at the pallet boundary are handled implicitly: if
\texttt{y\_row[g] = 1} and \texttt{plate[g] = 0}, the antecedent demands
\texttt{y\_row[g2] = 0}, a value outside the domain, so the implication is
trivially true. The plate opens onto empty space, as expected.

The formulation is correct but expensive for the propagator. The
disjunction on the plate value forces a reification on both branches; the
implication on the temporal precedence \(g_2 > g\) introduces a further
reified guard on every pair.

\subsection{Refined formulation}

The refined formulation removes both reifications. First, it introduces an
auxiliary variable that materialises the target row of each plate:

\vspace{0.5cm}
\begin{lstlisting}
array[1..n_chunks] of var 0..boxes_along_y+1: target;
constraint forall(g in 1..n_chunks)(
  target[g] = y_row[g] - 1 + 2*plate[g]
);
\end{lstlisting}
\vspace{0.5cm}

When \texttt{plate[g] = 0} the expression yields \texttt{y\_row[g] - 1};
when \texttt{plate[g] = 1} it yields \texttt{y\_row[g] + 1}. The two
opening directions collapse into a single arithmetic equality, replacing
the reified disjunction of the initial formulation. The domain is widened
to \([0, \texttt{boxes\_along\_y}+1]\) so that the boundary values \(0\)
and \(\texttt{boxes\_along\_y}+1\) explicitly represent "off the pallet";
they never match a real row index and so they encode a free edge without
any special case.

Second, the temporal precedence is built directly into the quantifier
range. Because the position of a chunk in the belt stream coincides with
its release index, only pairs with \(g_2 < g\) describe a potential
collision (\(g_2\) is already on the pallet when \(g\) is released). The
range \texttt{g2 in 1..g-1} is a compile-time parameter, so the solver
generates only the relevant pairs:

\vspace{0.5cm}
\begin{lstlisting}
constraint forall(g in 2..n_chunks, g2 in 1..g-1)(
  not(
       y_row[g2] = target[g]
    /\ x_start[g2] < x_start[g]  + chunk_len[g]
    /\ x_start[g]  < x_start[g2] + chunk_len[g2]
  )
);
\end{lstlisting}
\vspace{0.5cm}

The body becomes a direct negation of a single conjunction: chunk \(g_2\)
cannot share the target row of \(g\) while overlapping it along \(x\). No
reification on the plate value (encoded by \texttt{target}), no
reification on the temporal ordering (encoded by the quantifier range).
A single propagator per pair, where the previous formulation needed
several.